\section{Comparison: Standard ReCom, Merge-Split, Forest ReCom}
\label{sec:comparison}

\subsection{Three-Method Landscape}

Table~\ref{tab:comparison} summarises the key properties of the three
main reversible-ReCom designs.
Standard \recom{}~\citep{deford2021} is included as the baseline; it
achieves high throughput but leaves the finite-chain target implicit.
\mergesplit{}~\citep{janson2022} adds an explicit two-tree correction for
balanced-cut-count asymmetry.
\forestrecom{}~\citep{mccartan2020,dellaLibera2026forestRecom} supplies the
stronger reversible-ReCom reference point, with implementation details and
costs depending on whether one uses determinant-style forest counts or a
UST-family approximation.

\begin{table}[ht]
\centering
\caption{Comparison of the three main ReCom-family MCMC chains.
  Acceptance rates and EC quality are empirical estimates from our
  experiments on NC ($k=14$) and WI ($k=8$); all values are approximate
  and state-dependent.
  $m$ = edges in merged region $\approx 2n/k$ on average.}
\label{tab:comparison}
\renewcommand{\arraystretch}{1.35}
\small
\begin{tabular}{@{}p{0.22\linewidth}p{0.22\linewidth}p{0.24\linewidth}p{0.24\linewidth}@{}}
\toprule
Property & Standard \recom{} & \mergesplit{} & \forestrecom{} \\
\midrule
Target claim
  & Implicit/unknown
  & Approx.\ two-tree correction
  & Stronger reversible forest model \\
Acceptance ratio
  & 1 always
  & cuts\_fwd/cuts\_rev
  & forest-count or UST-family ratio \\
Ratio type
  & N/A
  & Stochastic per-step proxy
  & Determinant or sampled proxy \\
Per-step cost
  & $O(m \log m)$
  & $O(m \log m)$
  & variant-dependent \\
Empirical accept rate
  & $\approx\!80\%$
  & $\approx\!60\%$
  & $\approx\!40\%$ \\
EC quality (NC/WI)
  & Moderate
  & High
  & High \\
\bottomrule
\end{tabular}
\end{table}

\begin{remark}[Distinct target distributions]
Forest \recom{} and \mergesplit{} should not be collapsed into the same target
claim.
\mergesplit{} exposes a two-tree cut-count correction and omits pair-count and
determinant terms.
Forest \recom{} uses a forest-based Metropolis correction and is the stronger
reversibility reference in this family.
For short practical runs the two may produce similar compactness and seat-share
profiles, but that is an empirical observation, not an identity of target
distributions.
\end{remark}

\subsection{Standard ReCom vs.\ Merge-Split}

Standard \recom{} achieves higher acceptance rates ($\approx\!80\%$ vs.\
$\approx\!60\%$) because it never rejects on distributional grounds---it
only rejects when no balanced cut exists in the forward tree.
\mergesplit{} additionally rejects when
$\mathrm{cuts\_fwd} / \mathrm{cuts\_rev} < 1$, which introduces a downward
penalty when the forward tree happened to have many balanced cuts (diluting
the proposal probability) while the reverse tree has few.

The acceptance rate penalty of \mergesplit{} relative to standard \recom{}
is the direct price of the declared stochastic correction.
In our NC experiments, this penalty is approximately~20 percentage points.
For a fixed computational budget, \mergesplit{} produces approximately 25\%
fewer accepted plans per second than standard \recom{}, but those plans are
drawn from a chain whose main cut-count correction is explicit rather than
from the unknown standard-ReCom stationary distribution.

The per-step cost is asymptotically equal: both run two Wilson's-algorithm
calls (standard \recom{} runs one, \mergesplit{} runs two, but the constant
factor is small).
In practice, the wall-clock overhead of the second Wilson's call is
approximately 1.5--2.0$\times$ per step, which is partially offset by the
lower acceptance rate (fewer plan updates per step).

\subsection{Merge-Split vs.\ Forest ReCom}

Both \mergesplit{} and \forestrecom{} make the acceptance ratio visible.
The key difference is what ratio is being estimated or computed:

\begin{itemize}
  \item Determinant-style \forestrecom{} computes a forest-count ratio using
        matrix-tree determinants~\citep{mccartan2020}.
        That gives a stronger formal target model at higher cost.
        UST-family \forestrecom{} implementations trade some of that exactness
        language for the same graph-sampling machinery used elsewhere in this
        crate family.

  \item \mergesplit{} estimates the ratio by sampling a second spanning
        tree.
        The estimate has variance and does not include the pair-count term.
        It is therefore best used as a fast, reproducible correction rather
        than as a stand-alone proof of exact uniform sampling.
\end{itemize}

In our empirical experiments on NC and WI, both chains produce similar
distributional profiles: the fraction of 7D/7R vs.\ 8D/6R plans in a
1{,}000-step ensemble differs by less than~2 percentage points between
\mergesplit{} and \forestrecom{}.
The practical difference is operational: \mergesplit{} has a compact
two-tree step record and the current CLI percentile selector, while
\forestrecom{} carries the stronger forest-measure lineage and should be chosen
when that claim is the point of the run.

\subsection{When Does the Exact Ratio Matter?}

The stronger \forestrecom{} ratio is valuable when:
\begin{itemize}
  \item The ensemble is small (few steps), and variance in the ratio
        estimate could systematically bias the distribution.
  \item The application requires stronger formal statements about the target
        distribution than the \mergesplit{} two-tree proxy supports.
  \item The merged-region graphs are small enough that richer forest-count
        accounting is affordable and the additional precision is worth the
        engineering cost.
\end{itemize}

For standard 50-state production search runs with $n \gtrsim 1{,}000$ tracts,
\mergesplit{} is preferred when the immediate output is a compact, reproducible
candidate-selection run rather than a formal forest-measure study.

\subsection{Empirical Results: NC and WI}

\paragraph{North Carolina ($k = 14$, $n \approx 2{,}200$).}
In our experiments with 1{,}000 steps from a \cs{}-optimal starting plan:
\begin{itemize}
  \item \mergesplit{}: acceptance rate $\approx\!62\%$, 620 accepted plans,
        ensemble EC range 2{,}180--2{,}310.
  \item \forestrecom{}: acceptance rate $\approx\!41\%$, 410 accepted plans,
        ensemble EC range 2{,}175--2{,}305.
  \item Standard \recom{}: acceptance rate $\approx\!79\%$, 790 accepted plans,
        ensemble EC range 2{,}150--2{,}380.
\end{itemize}
\mergesplit{} produces a tighter EC range than standard \recom{} (the
distributional constraint concentrates mass near the typical-cut region)
and comparable range to \forestrecom{}.
The seat distribution (fraction of plans with 7D/7R majority) is
$51\%$, $53\%$, and $49\%$ respectively---within sampling noise.
Across the 5 seeds, the fraction of plans at $7D/7R$ ranged from~48\% to~54\%,
giving an approximate standard error of $\pm 3\%$ --- consistent with the
4 percentage-point difference across methods being within sampling variation.

We ran 5 independent chains ($s \in \{42, \ldots, 46\}$) of $T=1{,}000$ steps
on NC and computed $\hat{R} < 1.12$ for Democratic seat share across chains,
indicating approximate convergence of the marginal seat distribution (though
$1{,}000$ steps is likely too short for full plan-space convergence).
Acceptance rates showed less than 3 percentage-point variation across seeds.

\paragraph{Wisconsin ($k = 8$, $n \approx 1{,}400$).}
In our experiments with 1{,}000 steps:
\begin{itemize}
  \item \mergesplit{}: acceptance rate $\approx\!58\%$, seat distribution
        $3D/5R$: $38\%$, $4D/4R$: $47\%$, $5D/3R$: $15\%$.
  \item \forestrecom{}: acceptance rate $\approx\!38\%$, seat distribution
        $3D/5R$: $40\%$, $4D/4R$: $45\%$, $5D/3R$: $15\%$.
  \item Standard \recom{}: acceptance rate $\approx\!81\%$, seat distribution
        $3D/5R$: $34\%$, $4D/4R$: $50\%$, $5D/3R$: $16\%$.
\end{itemize}
Again, \mergesplit{} and \forestrecom{} produce similar seat distributions,
while standard \recom{} shows slight differences attributable to the
unknown-distribution bias.
